\section{Limitations}\label{sec:limitations}

\paragraph{Calibration and labels.} The controller needs an operator label for the traffic domain and a per-domain calibration of the frontier and the learning curve. On an unlabelled or uncalibrated domain it does not train; unlabelled production traffic needs online estimates of both, which we leave open.

\paragraph{Scope of the live evaluation.} We evaluate two engines, vLLM 0.27.1 and, with five of the \nv{n.pol} policies, SGLang 0.5.10, one GPU type (H100), two targets of one model family, EAGLE-3 drafters with $k{=}3$, greedy decoding with outputs of at most 512 tokens that end naturally (with every output forced to 256 tokens, the adapted code drafter reaches only \nv{proto.fixed.n3200}$\times$ plain decoding; \cref{tab:protocol}), single-GPU serving, and prefix caching disabled. Drafter hot swap was tested functionally with tensor parallelism before the freeze, but the controller was not evaluated there. \pol{TIDE} reimplements only TIDE's training triggers, on a shared GPU and with a live drafter comparison; TIDE itself trains on separate GPUs. The live workload is a synthetic chat--new-domain--chat sequence; other traffic shapes are covered only in simulation.

\paragraph{Statistics.} Each cell has three paired repetitions, and runs were spread over \nv{hosts} nodes of the same type. Node interference slowed \nv{scan.flagged.word} baseline runs, which we replaced by reruns (\cref{app:rerun}); smaller slowdowns remain in the data. Paired differences of a few percent between adaptive policies, such as \pol{Full} versus \pol{Informed} on Qwen3-8B code, are within the run-to-run spread.

\paragraph{Development and evaluation.} The rules and constants were developed on runs that drew from the same prompt pools as the frozen evaluation, in different orders. The frozen evaluation therefore does not test the constants on unseen prompts; the simulator validation on unseen traces only partly fills this gap, because it shares the calibration.
